% MAEG4060 — User Guide (standalone PDF)
\documentclass[11pt,a4paper]{article}

\usepackage[margin=2.5cm]{geometry}
\usepackage[T1]{fontenc}
\usepackage{lmodern}
\usepackage{microtype}
\usepackage{booktabs}
\usepackage{url}
\usepackage[hidelinks,breaklinks]{hyperref}

\makeatletter
\g@addto@macro\UrlBreaks{\do\/\do-\do\.\do\=}
\makeatother

\title{%
  MAEG4060 Course Project\\
  \large User Guide}
\author{%
  Interactive Physical Simulation of Off-Road Vehicle Dynamics\\
  \small Guo Tsun Hei; Yim Fu Chong}
\date{April 2026}

\begin{document}
\maketitle

\section{How to start the application}

\begin{itemize}
  \item Open the \texttt{Game\_Build} folder in your submission package (or run from the repository root after a local build; see below).
  \item \textbf{Double-click} \texttt{maeg4060\_stage2.exe} in that folder to launch the program.
  \item Keep the executable beside the \texttt{Assets} folder so paths such as \texttt{Assets/models/model\_jeep.x} resolve. If models fail to load, confirm the working directory contains \texttt{Assets}.
\end{itemize}

\subsection{Building from source (optional)}

The project is a CMake target (\texttt{maeg4060\_stage2}) that needs legacy \textbf{D3DX} headers and import libraries (\texttt{d3dx9.h}, \texttt{d3dx9math.h}, etc.), which are \emph{not} supplied by the Windows 10/11 SDK alone. CMake resolves them in one of two ways: (1)~set \texttt{DXSDK\_DIR} to the root of the \textbf{June 2010 DirectX SDK} (must contain \texttt{Include} and \texttt{Lib}); or (2)~leave \texttt{DXSDK\_DIR} unset and let CMake download the \texttt{Microsoft.DXSDK.D3DX} NuGet package into the build directory on first configure (requires network access; configure may take several minutes).

\textbf{Recommended (matches \texttt{scripts/configure\_build\_release.bat}):} run \texttt{vcvars64.bat} from your Visual Studio installation so \texttt{cl} is on \texttt{PATH}, then from the repository root use the Ninja generator and select MSVC explicitly, for example:
\begin{verbatim}
cmake -B build -G Ninja -DCMAKE_BUILD_TYPE=Release ^
  -DCMAKE_CXX_COMPILER=cl -DCMAKE_C_COMPILER=cl
cmake --build build
\end{verbatim}
On some setups the Visual Studio CMake generator fails to detect a C++ compiler even after \texttt{vcvars64}; Ninja with \texttt{-DCMAKE\_CXX\_COMPILER=cl} avoids that. If you change the generator in an existing build tree, delete the \texttt{build} folder (or remove \texttt{CMakeCache.txt} and the \texttt{CMakeFiles} directory) before re-running \texttt{cmake}.

If the compiler stops with \textbf{C1060} (out of heap space), retry the build with less parallelism, e.g.\ \texttt{cmake --build build -j 1}, and close other heavy applications.

Alternatively, from the repository root: \texttt{cmake -B build -G "Visual Studio 17 2022" -A x64}, then \texttt{cmake --build build --config Release}, provided the toolset is detected correctly.

To package binaries: run \texttt{scripts/package\_game\_build.bat} to populate \texttt{Game\_Build}, or \texttt{cmake --install build --prefix Game\_Build}.

\subsection{GitHub Actions (optional)}

If \texttt{.github/workflows/windows-build.yml} is enabled on your fork, the \textbf{Windows build} workflow on \texttt{windows-latest} produces an artifact (zip) containing \texttt{maeg4060\_stage2.exe}, \texttt{D3DX9\_43.dll} when the NuGet D3DX path is used, and a copy of \texttt{Assets/}. Download the artifact, unzip on Windows, and run the executable from that folder (keep \texttt{Assets} beside the \texttt{.exe}).

\section{Hardware and software requirements}

\begin{itemize}
  \item \textbf{Platform:} Windows PC (64-bit Windows recommended).
  \item \textbf{Graphics API:} The application uses \textbf{DirectX~9.0c}. You need a DirectX~9-compatible GPU and the \textbf{DirectX~9 End-User Runtime} (or equivalent system components) installed as required by your Windows version.
  \item \textbf{Display:} A desktop resolution of at least \textbf{1280$\times$720} is recommended; other resolutions are supported when the window is resized.
  \item \textbf{Optional input:} An \textbf{XInput-compatible gamepad} (e.g.\ Xbox controller) for analog control. Keyboard control works without a gamepad.
\end{itemize}

\section{Controls}

\subsection{Keyboard (WASD and related keys)}

Input is handled via the Win32 keyboard API (\texttt{GetAsyncKeyState}) with a vehicle-style layout, for example:

\begin{center}
\begin{tabular}{@{}ll@{}}
  \toprule
  \textbf{Input} & \textbf{Action} \\
  \midrule
  W / S & Forward / reverse (throttle) \\
  A / D & Steer left / right \\
  Space & Brake (or hand brake, depending on build) \\
  M & Pause / resume windmill motion (Maya animation when present, or simple spin for the static mesh fallback) \\
  Esc & Quit application \\
  \bottomrule
\end{tabular}
\end{center}

\noindent
Exact key bindings may match the on-screen HUD or readme in the binary package; the layout above reflects the project design (WASD movement with separate brake).

\subsection{Gamepad (XInput)}

With an XInput controller connected, the application maps analog axes and triggers for smoother control than the keyboard:

\begin{center}
\begin{tabular}{@{}ll@{}}
  \toprule
  \textbf{Control} & \textbf{Action} \\
  \midrule
  Left stick & Steering \\
  Right trigger & Throttle (accelerate) \\
  Left trigger & Brake / partial reverse \\
  \bottomrule
\end{tabular}
\end{center}

\noindent
Triggers provide analog throttle similar to a real vehicle; keyboard W/S act as digital throttle inputs.

\section{Troubleshooting}

\begin{itemize}
  \item \textbf{Missing DLLs or DirectX errors:} Install or repair the DirectX~9 end-user runtime and ensure GPU drivers are up to date.
  \item \textbf{Black screen or failed device creation:} Try windowed mode, update graphics drivers, and confirm no other application has exclusive full-screen access to the GPU.
  \item \textbf{Assets not found:} Run from \texttt{Game\_Build} with the supplied \texttt{Assets} (or equivalent) folder beside the executable. For static prop replacements, place OBJ/MTL files under \texttt{Assets/models/replacements} (see \texttt{Assets/models/replacements/README.md}).
  \item \textbf{Gamepad not detected:} Use a wired connection or official wireless adapter; confirm the controller appears in Windows as an XInput device.
\end{itemize}

\end{document}
